% =========================================================
% Proof: Division Algorithm
% Source: volume-iii/algebra/abstract-algebra/notes/notes-integers.tex
% Reference: Aluffi, Algebra: Chapter 0
% =========================================================

\subsection*{Division Algorithm}
\label{prf:division-algorithm}

\begin{remark*}[Return]
\hyperref[thm:division-algorithm]{$\leftarrow$ Back to Theorem (Division Algorithm) in Notes}
\end{remark*}

\begin{proof}
\Claim Let $a, b \in \mathbb{Z}$ with $b \neq 0$. Then there exist unique
$q, r \in \mathbb{Z}$ such that $a = bq + r$ with $0 \leq r < |b|$.

\Given $a, b \in \mathbb{Z}$ with $b \neq 0$.
\Goal To find unique $q, r \in \mathbb{Z}$ with $a = bq + r$ and $0 \leq r < |b|$.
\Strategy We establish existence via the Well-Ordering Principle, then uniqueness
by a subtraction argument.

\medskip
\noindent\textit{Existence.}

\medskip
\Let $S := \{\, a - bx \mid x \in \mathbb{Z},\; a - bx \geq 0 \,\}$.
The set $S$ is nonempty (one can always choose $x$ sufficiently negative to make
$a - bx > 0$).
\ByThm{Well-Ordering Principle}, $S$ has a least element; call it $r$.

\Let $x$ be the integer achieving the minimum, so that
\begin{equation}
  a - bx = r. \tag{1}
\end{equation}
Since $r \in S$, we have $r \geq 0$. Moreover $r < |b|$: if $r \geq |b|$ then
$a - b(x \pm 1) = r \mp b$ would be a non-negative element of $S$ strictly
smaller than $r$, contradicting minimality.
\Therefore $x$ and $r$ satisfy $a = bx + r$ with $0 \leq r < |b|$.
Set $q := x$.

\medskip
\noindent\textit{Uniqueness.}

\medskip
Suppose $q', r' \in \mathbb{Z}$ also satisfy
\begin{equation}
  a - bq' = r', \qquad 0 \leq r' < |b|. \tag{2}
\end{equation}
Subtracting (2) from (1):
\[
  (a - bq) - (a - bq') = r - r'
  \;\implies\;
  0 = b(q' - q) + (r' - r). \tag{3}
\]
\Therefore $b(q' - q) = r - r'$, so $|b|\,|q' - q| = |r' - r|$.

Since $0 \leq r, r' < |b|$, we have $|r' - r| < |b|$.
If $q' \neq q$ then $|b|\,|q' - q| \geq |b|$, contradicting $|r' - r| < |b|$.
\Therefore $q' = q$, and (3) then gives $r' = r$.

\Hence $q$ and $r$ are unique. \AsReq
\end{proof}

\begin{remark*}[Proof shape]
Two-part direct proof. Existence uses the Well-Ordering Principle on the set
$S = \{a - bx \geq 0 \mid x \in \mathbb{Z}\}$: its least element is the
remainder $r$, and minimality forces $r < |b|$.
Uniqueness is an arithmetic squeeze: subtracting two representations gives
$b(q'-q) = r-r'$, and the bound $|r'-r| < |b|$ forces $q' = q$ and then $r' = r$.
This is the standard pattern for uniqueness in divisibility arguments —
assume two solutions, subtract, use the size bound to eliminate the non-trivial case.

\FlashProof{Two-part direct proof. Existence uses the Well-Ordering Principle on the set $S = \{a - bx \geq 0 \mid x \in \mathbb{Z}\}$: its least element is the remainder $r$, and minimality forces $r < |b|$. Uniqueness is an arithmetic squeeze: subtracting two representations gives $b(q'-q) = r-r'$, and the bound $|r'-r| < |b|$ forces $q' = q$ and then $r' = r$. This is the standard pattern for uniqueness in divisibility arguments — assume two solutions, subtract, use the size bound to eliminate the non-trivial case.}
\FlashUses{\begin{itemize} \item \hyperref[thm:wop]{Well-Ordering Principle} — to guarantee $S$ has a least element. \item \hyperref[def:divisibility]{Definition (Divisibility)} — context for what $b \mid a$ means when $r = 0$. \item Absolute value and integer arithmetic — $|b(q'-q)| = |b||q'-q|$ and the bound $|r'-r| < |b|$. \end{itemize}}
\end{remark*}

\begin{remark*}[Dependencies]
\begin{itemize}
  \item \hyperref[thm:wop]{Well-Ordering Principle} — to guarantee $S$ has a
        least element.
  \item \hyperref[def:divisibility]{Definition (Divisibility)} — context for
        what $b \mid a$ means when $r = 0$.
  \item Absolute value and integer arithmetic — $|b(q'-q)| = |b||q'-q|$ and
        the bound $|r'-r| < |b|$.
\end{itemize}
\end{remark*}